\section{Experimental Protocols}
\label{app:protocols}
\label{app:protocol}
FIT determines coordinates, gates and adapters; VAL sets prescribed movement
endpoints; HEAD trains independent readers; EVAL supplies participant scores.
Participant roles are disjoint within a split. Table~\ref{tab:protocol-map}
distinguishes native outputs, matched movement and task-informed P300 DEV
selection. Downstream labels construct associations and train heads, but
enter adapter selection only in reuse and never adapter fitting.
The repository README (Appendix~\ref{app:reproduction}) indexes participant
lists, executable configurations and per-study reproduction instructions.

\begin{table}[htbp]
\centering\small
\caption{\textbf{Protocol map.} Participant counts are FIT/VAL/HEAD/EVAL
(DEV replaces VAL in reuse). Gates are fixed during correction.}
\label{tab:protocol-map}
\begin{tabular}{@{}>{\raggedright\arraybackslash}p{.18\linewidth}>{\raggedright\arraybackslash}p{.43\linewidth}>{\raggedright\arraybackslash}p{.33\linewidth}@{}}
\toprule
Study & Coordinates, calibration and gate & Participants and operating point \\
\midrule
Controlled SSVEP & Pooled EEGPT512; full same-trial CAR/Oz contrasts;
90\% directional energy, ranks $3,3,3,2$ & Four role rotations;
81 distinct EVAL participants. Native correction/target; $.05$ VAL budget for permission. \\
\addlinespace
Temporal N170 & 240 temporal features; simultaneous Neuroscan/Flex pairs;
90\% energy, rank 56 & $4/2/4/10$; native maps and matched-budget IGBP comparison. \\
\addlinespace
Recorded SSVEP & Mean512 correction, $15\times512$ readouts;
six virtual FIT-local contrasts; 99\% energy, rank 328 & $41/20/20/21$;
native maps. \\
\addlinespace
P300-selected reuse & EEGPT/PCA256, P300 FIT reference coordinates;
400 simultaneous CAL pairs, rank 8 & $4/2/4/8$;
P300 DEV-selected map, unchanged on N170/MMN. \\
\bottomrule
\end{tabular}
\end{table}

\subsection{Measurement views and participant roles}
\paragraph{Controlled SSVEP.}
The wet/dry collection contains 102 participants, 12 frequencies and ten
blocks per frequency and electrode type \citep{zhu2021ssvep}. Wet trials
from eight occipital channels are cropped to samples 160--659 at 250 Hz,
standardized within channel and encoded by frozen EEGPT
\citep{wang2024eegpt}. CAR and Oz-reference views are encoded separately;
averaging 15 tokens yields $h_i^{\rm CAR},h_i^{\rm Oz}\in\R^{512}$.
The evaluated endpoints are
\[
 h_i(u)=\tfrac12(h_i^{\rm CAR}+h_i^{\rm Oz})
 +(2u-1)\tfrac{\gamma}{2}(h_i^{\rm Oz}-h_i^{\rm CAR}),\qquad\gamma=.005,
\]
whereas calibration uses the full CAR--Oz contrasts. The task is frequency
$\geq12$ Hz. Within each FIT/VAL participant--frequency cell, nine of ten
trials receive source 1 in the high-frequency class and one of ten in the
low-frequency class. Adapters fit these observed endpoints; HEAD/EVAL
readers use both endpoints with equal source mass.

Five folds are stratified by electrode-session order and gender. In rotation
$r$, folds $r,r+1,r+2$ supply EVAL, HEAD, VAL cyclically, and the remaining
two supply FIT. The four FIT/VAL/HEAD/EVAL counts are $41/20/20/21$,
$42/20/20/20$, $41/21/20/20$ and $40/21/21/20$. Each participant supplies
120 underlying trials. The 81 distinct EVAL participants appear once per
fit seed; source assignment is fixed across seeds. Closed-form comparisons
use one deterministic map per split and the seed-17 reader realization;
gate/target comparisons cross these rotations with all shared fit seeds.

\paragraph{Temporal N170.}
Neuroscan (source 0) and Flex (source 1) record the same watch/face events
concurrently \citep{williams2020flex}. Processing uses 16 common channels,
0.1--30 Hz zero-phase FIR filtering, CAR, resampling to 128 Hz, and one-second
epochs starting at $-.2$ s with prestimulus baseline subtraction. Event
ordinals and timestamps align the views. A pair is retained only if neither
device exceeds $100\,\mu$V absolute amplitude after baseline correction.
FIT/VAL/HEAD/EVAL contain $622/288/716/1381$ clean pairs; every EVAL
participant retains both classes. Fifteen successive 50 ms channel averages
from 0 to 750 ms give 240 coordinates without an encoder or PCA. This is
distinct from the EEGPT-based reuse protocol below.

\paragraph{Recorded SSVEP.}
Source 0/1 denotes actual dry/wet sessions, with no physiological trial
correspondence. The roles contain $9840/4800/4800/5040$ observations under
one fixed stratified split. Six calibration transformations of each trial
use spectral transport, spatial transport and their composition at strengths
.5 and 1 toward the opposite source. FIT prestimulus signals determine
smoothed log-power gains clipped to $[.25,4]$ and spatial whitening/coloring
with covariance shrinkage .1. MGPA corrects the token mean and adds its
displacement to every token; source and twelve-frequency task readers
access the full $15\times512$ representation.

\subsection{Fitting, replay and comparison interfaces}
\label{app:fitting-replay}
\label{app:component-fitting}
\begin{table}[htbp]
\centering\small
\caption{\textbf{Shared MGPA settings.} Component ablations change only the
specified critic or target. Split, reader and assignment seeds are separate.}
\label{tab:shared-settings}
\begin{tabular}{@{}>{\raggedright\arraybackslash}p{.34\linewidth}>{\raggedright\arraybackslash}p{.60\linewidth}@{}}
\toprule
Component & Setting \\
\midrule
Iterative fits / horizon & Seeds 17, 29, 43; at most 48 stages \\
Source critics & Balanced logistic ($C=1$); GELU MLP, hidden width 128 \\
Conditional-mean anchor & Ridge with intercept and penalty 10 (both constructions) \\
Neural-critic optimizer & AdamW, batch 256, weight decay $.001$ \\
Gradient stabilizer / damping & $.1$ / $.001$ \\
Controlled and real-representation fits & 12 critic epochs, LR $.001$, radius $.05\sqrt d$ \\
P300 reuse fits & At most 100 epochs, LR $.0003$, source-VAL patience 10,
radius $.025\sqrt d$ \\
Closed-form residual covariance & $.01$ shrinkage toward mean-diagonal identity \\
Participant uncertainty & 20,000 percentile bootstrap resamples, seed 1101 \\
\bottomrule
\end{tabular}
\end{table}

\begin{algorithm}[htbp]
\small
\caption{Iterative MGPA: fit once, replay without source labels}
\label{fig:algorithm-main}
\begin{enumerate}
\setlength{\itemsep}{2pt}
\item \textbf{Inputs:} calibration pairs, source-labeled FIT/VAL embeddings,
gate rank or energy threshold, critic bank, $\lambda,\tau,\rho,K$, and
endpoint rule. Task-informed selection additionally requires DEV task data.
\item Fit standardization and the contrast-SVD gate \eqref{eq:gate-estimate};
freeze $B,P,Q$. Return identity if $r=0$. Store original FIT/VAL complements.
\item Fit source critics on current FIT embeddings and regress their scores
on original FIT $Q$ to obtain conditional anchors.
\item Store critics and anchors; compute $e,A,W$, solve
\eqref{eq:tr-problem}, and apply $h\leftarrow h+B\Delta z_*$. Record VAL
outputs and repeat Steps 3--4 up to $K$ stages.
\item Select and store the prefix and optional final-stage fraction using
the declared validation rule; only task-informed selection uses task labels.
\item \textbf{Deployment:} standardize one embedding, replay the stored map,
and return it in the downstream interface's coordinates.
\end{enumerate}
\end{algorithm}

Standardization uses FIT only. Gates use row-normalized contrasts;
controlled and N170 fits filter contrast norms at the fifth percentile,
whereas recorded SSVEP uses no percentile filter. Controlled gating uses
at most 2048 deterministically sampled contrasts. Energy thresholds and
ranks appear in Table~\ref{tab:protocol-map}.
Closed-form MGPA uses the same gate and standardization, the mean projected
calibration contrast, and the shared Ridge anchor. Its metric comes from
residuals of a joint regression of $Z$ on $Q$ and source: intercept/source
coefficients are unpenalized and $Q$ has penalty 10. The residual second
moment is normalized by FIT count and shrunk as in
Table~\ref{tab:shared-settings}; \eqref{eq:closed-form-mgpa} uses linear
solves. Native studies take $\alpha=1$; reuse selects $\alpha$ on DEV.
Recorded contrasts are oriented wet-like to dry-like for both gate and direction.

The target ablation replaces the conditional anchor by the current-stage
empirical FIT-mean score; a zero-score control is also reported. The
linear-only ablation removes the nonlinear critic. Other fitting settings
remain matched as specified in the corresponding results.

\paragraph{Baselines and deployment information.}
LEACE fits an unrestricted covariance-weighted affine eraser to observed
FIT rows. CORAL uses FIT source means and identity-regularized sample
covariances, not pair correspondence. FEATMAP fits centered unregularized
least squares to FIT pairs, with the minimum-norm slope for rank-deficient
designs; controlled pairs use the evaluated $\gamma$-scaled endpoints.
CORAL/FEATMAP leave their chosen reference unchanged and require source
routing at deployment; both orientations are reported. MGPA, LEACE and
IGBP deploy source-blind. Recorded SSVEP has no paired FEATMAP comparison.
No method uses EVAL moments. Reuse searches regularization/strength on
P300 DEV instead of using these native settings.

\paragraph{IGBP.}
\label{app:igbp-protocol}
IGBP~\citep{iskander2023} shares each study's roles, standardized inputs and
independent readers. At each iteration a $d\to d\to2$ ReLU classifier
drives an ungated projective update without a conditional anchor or pair
correspondence. All reported fits use the longer-training preset: AdamW,
LR $2\times10^{-4}$, weight decay $.01$, batch 256, at most 50 epochs,
VAL source-accuracy patience 10 and improvement threshold $.002$, with no
$.80$ accuracy exit. Native maps apply 100 projections without task-based
stopping. Updates are evaluated directly in logit space to avoid softmax
saturation. For N170, $d=240$; for recorded SSVEP, $d=512$ and the resulting
mean correction is added to every token, matching MGPA's interface.

\paragraph{Matched movement and edit-permission controls.}
\label{app:gate-attribution}
Native outputs apply the complete fitted map; movements need not match.
A budget $b$ instead selects the first
upward crossing of $\|H'_{\rm VAL}-H_{\rm VAL}\|_F/\|H_{\rm VAL}\|_F=b$,
interpolating only the last update and replaying that prefix/fraction on
HEAD/EVAL. Controlled permission compares the measurement gate, leading
observed-FIT PCA directions of the same rank, five fixed random orientations
of that rank (averaged, not selected), and all 512 directions. Each arm
refits its critics and anchors with the same recipe and original measurement
complement $Q_0$ as anchor input. Thus controls receive calibration-derived
rank and conditioning information; only the measurement arm necessarily
preserves $Q_0$. Every correction arm and random orientation attains the
reported $.05$ budget; all 108 cells, including identity, are available.
At the originally primary $.10$ budget, 17 of 60 random trajectories
do not attain it within the horizon; native endpoints are not substituted.
The N170 IGBP budget comparison uses the same first-crossing rule.

\subsection{Independent readers and statistical summaries}
\label{app:independent-readouts}
\label{app:full-token-readers}
Source readers fit corrected HEAD outputs independently of adapters and
score participant-disjoint EVAL outputs. Table~\ref{tab:reader-bank} specifies
their inputs and capacity. All preprocessing uses HEAD only; each recording,
not each token, is one example. Neural readers train without early stopping
and contribute three restarts at 12/50/200 epochs.

\begin{table}[htbp]
\centering\small
\caption{\textbf{Independent source-reader bank.} Controlled/N170 use the
11 common readers. Recorded SSVEP uses all $11+20+10=41$ readers; its common
bank alone defines $S_{\rm mean}$.}
\label{tab:reader-bank}
\begin{tabular}{@{}>{\raggedright\arraybackslash}p{.23\linewidth}>{\raggedright\arraybackslash}p{.71\linewidth}@{}}
\toprule
Input & Readers and training \\
\midrule
Common representation & Ordinary and covariance-whitened logistic ($C=1$,
whitening shrinkage $.01$); GELU MLP with width 128, AdamW, LR $.001$,
batch 256, weight decay $.001$. \\
\addlinespace
Direct full tokens & Raw logistic ($C=1$), HEAD-standardized logistic
($C=.01$); flat $7680\to128\to2$ GELU with common training; token network
with LayerNorm, shared $512\to64$ GELU, concatenation and binary output,
LR $.0005$, batch 128, weight decay $.001$, gradient clipping 5. \\
\addlinespace
Centered token residuals & Standardized logistic and the same token network
with its three restarts/checkpoints. \\
\bottomrule
\end{tabular}
\end{table}

Recorded common-offset correction preserves $R_t=H_t-\bar H$.
Residual readers fit once on original HEAD residuals and are reused on
corrected EVAL; numerical invariance is checked. This retained token
variation is distinct from the pooled gate complement $Q$.

\paragraph{Task utility and movement.}
Controlled/N170 heads are HEAD-standardized balanced Ridge classifiers
(penalty 10, LSQR). The recorded task head uses HEAD-standardized flattened
features and the token architecture above with a twelve-class output,
trained for 50 epochs. Frozen heads and scalers train on original HEAD;
refitted heads and scalers train on corrected HEAD; both evaluate corrected
EVAL. $T_f/T_r$ equally average participant AUROCs over $0\to1$ and
$1\to0$ transfer; recorded scores are twelve-class macro AUROC.
$M$ averages $\|H'_p-H_p\|_F/\|H_p\|_F$ over EVAL participants, which
can differ from aggregate VAL movement.

\paragraph{Aggregation and uncertainty.}
\label{app:evaluation-scope}
Source accessibility maximizes participant-mean AUROC over readers within
each fit, then averages fits. Controlled scores pool disjoint EVAL folds;
deterministic maps contribute one output per split. Shared participant
bootstrap draws pair methods, transfer directions and fits, recomputing
the source maximum inside every draw. Intervals condition on fitted
coordinates, adapters and readers and are pointwise, without multiplicity
adjustment; seeds/trials do not increase participant counts. These cohorts
informed protocol development, including the P300 selection task and rank
eight. Controlled gate/target studies share participants; the two N170
protocols overlap. Source maxima measure accessibility under the specified
bank, not complete erasure or unbiased performance of an EVAL-selected reader.

\subsection{P300-selected correction reused on N170 and MMN}
\label{app:recovery-protocol}
\label{app:reuse-protocol}
\paragraph{Calibration and shared coordinates.}
All three tasks use simultaneous Neuroscan/Flex recordings and the same
eight EVAL participants. The first 100 consecutive P300 pairs per FIT
participant form CAL. Any other FIT event whose one-second epoch overlaps
CAL on either device is excluded, leaving 2260 observed events. One endpoint
per observed event plus both endpoints of the 400 CAL pairs yields 3060
unique fitting endpoints for MGPA, LEACE and CORAL; FEATMAP uses the 400
pairs. CAL supplies both MGPA gates and the Closed-form direction;
Iterative critics fit the pooled bank without paired targets.

Frozen EEGPT retains $7\times4\times512=14{,}336$ features from seven
overlapping windows, with fixed $20\,\mu$V input scaling. Centered PCA256
and standardization fit only the Neuroscan endpoints of remaining P300 FIT
events, then apply to both devices and all tasks. Neither N170 nor MMN
fits coordinates or selects the adapter. HEAD uses all technically valid
events, unlike the amplitude-filtered temporal N170 study.

\paragraph{Frozen heads and association shifts.}
Each task has two balanced logistic heads trained on original HEAD features
with equal participant weights: an aligned, shortcut-sensitive head and an
independent-assignment control head. Regularization is selected from
$C\in\{10^{-6},10^{-5},\ldots,10\}$ by leave-one-HEAD-participant-out
independent-head AUROC and shared between that task's heads. Heads,
standardizers and assignment realizations stay fixed across adapter fits.
For positive-class prevalence $\pi_p$, endpoint assignment satisfies
\[
 \Pr(U=1\mid Y=y,p;\xi)=\tfrac12+.4\xi
 \frac{y-\pi_p}{\max(\pi_p,1-\pi_p)},\qquad \xi\in[-1,1].
\]
Aligned/independent/reversed settings are $\xi=1/0/-1$; device marginals
remain one half and signals are unchanged. Fitting samples recorded
endpoints; evaluation weights both endpoints. Worst AUROC is the exact
minimum of the participant-mean quadratic AUROC curve on $[-1,1]$,
including interior minima. Control-head AUROC is evaluated at $\xi=0$.

\paragraph{Selection and transfer.}
Both MGPA constructions use rank-eight gates and pooled conditional-mean
targets. Iterative training follows Table~\ref{tab:shared-settings}, with
1332 DEV source-validation endpoints. P300 DEV selects the endpoint
maximizing worst AUROC subject to at most $.01$ loss from identity for
\emph{each} head at independent association. Ties prefer less movement,
then an earlier endpoint. All iterative prefixes are eligible. Affine
strengths are $0,.02,\ldots,1,1.1,1.2,1.3$; baseline ridge grids and both
CORAL/FEATMAP reference choices are searched under the same rule. Across
recipes within a family, candidates within $.002$ of the best admissible
DEV worst AUROC are ordered by movement and the fixed tie rule. The grids
and selected values are recorded in the repository's task-reuse guide;
Closed-form MGPA selects $\alpha=1.1$. Each complete selected map is then
reused unchanged on P300, N170 and MMN.

Reuse movement for DEV tie-breaking is the participant mean of the ratio
of mean paired-event displacement norm to mean input norm, rather than the
Frobenius ratio above. Reported worst AUROC averages fit-specific minima
of participant-mean curves, not the minimum of an ensemble curve.
Each shared participant-bootstrap draw recomputes every curve and its
minimum before averaging fits; the same draws pair all tasks and methods.
Control-head AUROCs average participants and fits. Other uncertainty
conventions follow Appendix~\ref{app:evaluation-scope}.
